\section*{Exact four-channel law for a shared physical query}

All logarithms are natural.  For a real function $c$ on a finite cube put
\[
 \mathcal E(c)=\sum_U \widehat c(U)^2
      \log\frac1{\widehat c(U)^2},
\]
with the zero terms omitted, and define the polarized entropy kernel
\[
 \mathfrak D(p,q)=\mathcal E(p)+\mathcal E(q)
 -\frac{\mathcal E(p+q)+\mathcal E(p-q)}2.
\tag{1}
\]
Equivalently, $\mathfrak D(p,q)=\sum_U K(\widehat p(U),\widehat
q(U))$ for the scalar kernel used in the preceding B12 runs.

Let $g_+,g_-,h$ be Boolean signs on the same cube, set
\[
 a=(g_++g_-)/2,\qquad b=(g_+-g_-)/2,\qquad F=a+hb,
\]
and query a physical coordinate $i$ shared by all three functions.  Write
\[
 a=a_0+x_i a_1,\qquad b=b_0+x_i b_1,\qquad h=r+x_i s.
\]
On the cube with coordinate $i$ removed define
\[
 A=a_0+r b_0,\quad B=a_1+r b_1,\quad
 C=s b_1,\quad D=s b_0.                                      \tag{2}
\]
Then multiplication before restriction gives the exact split
\[
 F=(A+C)+x_i(B+D).                                           \tag{3}
\]

For a Boolean sign $f=u+x_i v$, its physical entropy decrement is
\[
 D_i(f)=\mathfrak D(u,v).
\]
Put $q=\|b\|_2^2$ and introduce the source channel, crossed channel and
polarized interaction
\[
 S_i=\mathfrak D(A,B),\qquad X_i=\mathfrak D(C,D),             \tag{4}
\]
\[
 \Omega_i=\mathfrak D(A+C,B+D)-S_i-X_i.                       \tag{5}
\]
Finally define the signed shared remainder
\[
 R_i^{\rm sh}=\Omega_i-qD_i(h)
              =\Omega_i-q\mathfrak D(r,s).                    \tag{6}
\]

\subsection*{Theorem (proved)}

For every dimension, every Boolean triple $(g_+,g_-,h)$, and every shared
physical coordinate $i$,
\[
 \boxed{D_i(F)=S_i+X_i+qD_i(h)+R_i^{\rm sh}.}                 \tag{7}
\]
The statement includes all coherent Fourier collisions; no disjointness,
orthogonality, bias floor, or collision separation is assumed.

\paragraph{Proof.}
Equation (3) gives $D_i(F)=\mathfrak D(A+C,B+D)$.  Substituting (5) and
then (6) proves (7).  Although the algebra is short, the placement of the
four channels is essential: products in (2) are taken in physical space, so
their Fourier coefficients are the actual signed convolutions on the
remaining shared coordinates.  No tensor identity has been used. \(\square\)

\subsection*{Three exact boundary checks}

If the query belongs only to a disjoint child block, then $a_1=b_1=0$.
Thus $S_i=X_i=0$, and (6) becomes
\[
 R_i^{\rm sh}=D_i(F)-qD_i(h)=R_\mu(\alpha),                  \tag{8}
\]
the complete biased coherent remainder from the earlier B12 kernel theorem,
including its $-qK(\mu,\alpha)$ correction.

If the selector does not depend on $i$, then $s=0$, so $C=D=0$,
$\Omega_i=0$, $D_i(h)=0$, and $R_i^{\rm sh}=0$.  The entire decrement is
the source channel $S_i$.  Thus (7) does not incorrectly charge a child
remainder at a source-only query.

If the selected output and both restrictions are constant, then
$D_i(F)=0$.  Equation (7) retains the exact signed cancellation
\[
 S_i+X_i+qD_i(h)+R_i^{\rm sh}=0.                              \tag{9}
\]
This is precisely the cancellation erased by the nonnegative two-term
potential in the preceding constant-output obstruction.

\subsection*{The signed remainder alone does not repair the old Bellman step}

The tempting next claim would be
\[
 R_i^{\rm sh}\le
 \Delta_i\{C_*q\mathbf I(h)+\operatorname{Var}(h)
                    \mathcal E(b)\},
 \qquad C_*=4+2\log2+2/e.                                    \tag{10}
\]
It is false even for a genuine activation-irreducible three-bit selector.
In zero-based lexicographic order take
\[
\begin{split}
g_+&=(-1,-1,1,1,-1,-1,1,-1),\\
g_-&=(-1,-1,-1,1,-1,-1,-1,-1),\\
h&=(-1,-1,1,-1,-1,1,-1,1),
\end{split}                                                   \tag{11}
\]
and query coordinate $i=1$.  The branches disagree exactly at inputs
$\{2,6\}$, where $h$ has signs $+1$ and $-1$ respectively.  Hence both
branches are actually selected on disagreement points.

Exact Walsh arithmetic yields
\[
\begin{aligned}
S_i&=\frac{15}{4}\log2-\frac{25}{32}\log5-\frac9{32}\log3,\\
X_i&=\frac14\log2,\\
D_i(h)&=3\log2-\frac98\log3,\qquad q=\frac14,\\
R_i^{\rm sh}&=-\frac{11}{4}\log2+\frac{25}{32}\log5
                         +\frac9{16}\log3,\\
D_i(F)&=2\log2.
\end{aligned}                                                 \tag{12}
\]
The four lines in (12) satisfy (7) exactly.  For the old potential,
\[
 \Delta_i(q\mathbf I(h))=-\frac1{16},\qquad
 \Delta_i(\operatorname{Var}(h)\mathcal E(b))=\frac7{16}\log2,
\tag{13}
\]
and therefore
\[
 \Delta_iP_*=-\frac14-\frac1{8e}+\frac5{16}\log2.            \tag{14}
\]

The separation is rigorous.  Since
\[
 163\,5^{25}3^{18}>60\,2^{88},\qquad
 e>\sum_{k=0}^5\frac1{k!}=\frac{163}{60},
\]
we obtain $5^{25}3^{18}/2^{88}>1/e$, and hence (12) gives
\[
 R_i^{\rm sh}>-\frac1{32}.                                   \tag{15}
\]
Also $\log2<7/10$, because the first five terms of the exponential series
at $7/10$ already exceed $2$.  Equation (14) therefore gives
\[
 \Delta_iP_*<-\frac14+\frac7{32}=-\frac1{32}.                 \tag{16}
\]
Combining (15)--(16) refutes (10).

This is an obstruction to separating $R_i^{\rm sh}$ from $S_i$ and $X_i$;
it is not a counterexample to FEI.  The exact local output decrement in this
example is positive, and (7) holds.  A surviving overlap argument must bundle
the source and crossed channels with the signed interaction across a block,
or define a future-activation gauge on which those channels telescope
together.  Enveloping the signed remainder alone by the old product-potential
drop loses an exact $0.048564\ldots$ nats at this one physical step.
